\scanpage{12}

\deleted{\unsure{puis schéma en groupes} $G$, de type fini sur $\mathbb Z$, qui \deleted{\unsure{repr.}} est représenté matriciellement par}

\unsure{$G$.} Considérons le schéma en groupes de $\operatorname{GL}(\unsure{n|3})_{\mathbb Z}$ qui correspond aux matrices de la forme
\diagram{\[\left(\begin{array}{ccc}1&\alpha&\gamma\\0&\lambda&\beta\\0&0&\mu\end{array}\right)\]}
(avec $\lambda\mu$ une unité bien sûr), \inserted{il opère sur $M$} de la façon évidente qui consiste à interpréter $ABCDEF$ comme coeff. d'une forme
\[
\deleted{\unsure{X^2}+AXY^2}\qquad X^2\unsure{+|-}AY^3+BXY^2+CY^{2}+DX+EY+F
\]
et en faisant sur $X,Y$, la substitution affine
\[
\unsure{Y|X}'=\lambda\unsure{Y|X}+\alpha,\quad X'=\mu X+\beta Y+\gamma
\]
[on peut aussi dire --- \unsure{ce sera plus clair}, que $G$ opère sur le \unsure{vecteur} $E_3$, donc dans $\Phi_3(E_3)$, et qu'il laisse invariant le sous-\deleted{module}\inserted{\unsure{schéma}} de \deleted{\unsure{dimension}}\unsure{rang} $6$ des formes qui sont de la forme dite, d'où les opérations de $G$ sur $M$.

\margin{\uuline{Prop}}
Mais le foncteur qui à tt $S$ attribue \unsure{l'ens.} \inserted{l'ens des classes à isom près} des courbes elliptiques ``structurées'' comme dessus sur $S$, est \uline{représentable} par $M$. De plus, pour que deux \unsure{éléments} de
\mspage{11}
